% !TeX spellcheck = de_DE
\documentclass[]{article}

\usepackage[utf8]{inputenc}
\usepackage[T1]{fontenc}
\usepackage[scaled]{beramono}
\usepackage{amssymb}
\usepackage{amsmath}
\usepackage{graphicx}
\usepackage[top=2cm, bottom=2cm, left=2cm, right=4cm]{geometry}
\usepackage{csquotes}
\usepackage{listings}
\usepackage{hyperref}
\usepackage{booktabs}

\lstset{
	tabsize=2,
	language=Java,
	showstringspaces=false,
	basicstyle=\footnotesize\ttfamily,
}

%no indentation
\setlength\parindent{0pt}

\pagestyle{empty}

\begin{document}

	\section{Programmieren auf Papier}
	
	\subsection{Felder}
	
	Schreiben Sie eine Funktion, die ermittelt, wie oft in einem Feld gleiche Werte aufeinander folgen.
	
	\begin{lstlisting}
		#include <stdio.h>
		
		int numpairs(int arr[], int length) {		
			int i, c = 0;
			for (i = 0; i < length - 1; i++) 
				if (arr[i] == arr[i + 1])
					c++;
			return c; 
		}
				
		int main() { 		
			int arr[] = {3,2,3,3,7,1,7,7,7};		
			int n = numpairs(arr, 9);		
			printf("%d\n", n); /* output: 3 */
			return 0;
		}
		
	\end{lstlisting}
	
	\subsection{Matrizen}
	
	Schreiben Sie eine Funktion, die 1 zurückgibt, wenn zwei Matrizen gleich sind, sonst 0.
	
	\begin{lstlisting}
		#include <stdio.h>
		#define ELEMENTS 4
		
		int equals(int first[ELEMENTS][ELEMENTS], int second[ELEMENTS][ELEMENTS]) {
			int i, j;
			for (i = 0; i < ELEMENTS; i++)
				for (j = 0; j < ELEMENTS; j++)
					if (first[i][j] != second[i][j])
						return 0;
			return 1; 
		}
				
		int main() {   
			int matrix1[ELEMENTS][ELEMENTS] = {
				{0,0,0,0},
				{1,0,0,0},
				{0,1,0,0},
				{0,0,1,0}};        
			int matrix2[ELEMENTS][ELEMENTS] = {
				{0,0,0,0},
				{1,0,0,0},
				{0,1,0,0},
				{0,0,1,0}};
			printf("%d\n", equals(matrix1, matrix2)); /* output: 1 */
			return 0;
		}
		
	\end{lstlisting}	
	
	\clearpage
	
	\section{Ausdrücke}
	
	Bestimmen Sie Typ und Wert der gefragten Ausdrücke \lstinline|???|.
	
	\begin{minipage}[t]{0.45\linewidth}
		\begin{lstlisting}[gobble=6]
			int i, j, x = 0;
			for(i = 0; i < 3; i++)
				for(j = 0; j < 3; j++)
					x++;
			x == 9 /* int */
		\end{lstlisting}
	\end{minipage}
	\quad
	\begin{minipage}[t]{0.55\linewidth}
		\begin{lstlisting}[gobble=6]
			int i, j, x;
			for(i = j = 0, x = 0; i < 3 && j < 3; i++ ,j++, x++);
			x == 3 /* int */
			
			
		\end{lstlisting}
	\end{minipage}
	%
	\begin{minipage}[t]{0.45\linewidth}
		\begin{lstlisting}[gobble=6]
			int a = 3;
			int b = 3;
			int *x = &a, *y = &b;
			(x == y) == 0   /* int */
			(*x == *y) == 1 /* int */
		\end{lstlisting}
	\end{minipage}
	\quad
	\begin{minipage}[t]{0.55\linewidth}
		\begin{lstlisting}[gobble=6]
			int a = 5, *x = &a;
			while(a > 1)
				(*x)--;
			a == 1 /* int */
			
		\end{lstlisting}
	\end{minipage}

	
	\section{Speicherverwaltung}
	
	Das folgende C-Programm wird auf einem Rechner mit 12 Speicherzellen ausgeführt, die jeweils eine Speicheradresse oder einen \lstinline|int|-Wert enthalten können.
	Neu deklarierte Variablen werden einer zufälligen freien Speicherzelle zugewiesen.
	Skizzieren Sie einen möglichen Zustand des Speichers nach jeder Programmzeile.\\
	
	\begin{minipage}[t][][b]{0.3\linewidth}
		\begin{lstlisting}[gobble=6,numbers=left]
			int i = 3, j[3] = {1, 5};
			int *x = &i;
			int *y = j + 1;
			*(y + 1) = 7;
			int **z = &y;
		\end{lstlisting}
	\end{minipage}	
	\qquad
	\begin{minipage}[t][][b]{0.65\linewidth}
		\begin{center}
			\includegraphics[width=\linewidth]{figures/memory-solution}		
		\end{center}
	\end{minipage}
	

\end{document}